% Transcription — fonds Alexandre Grothendieck, shelfmark 150, batch 4 (pages 61-80).
% Established from the facsimile archives/batches/150/batch-04.pdf (a copy of
% 150.pdf fetched through the project's relay,
% https://grothendieck-relay.onrender.com/source/150.pdf; PDF page k =
% archivists' page 60+k, no cover sheet), per the skill transcribe-grothendieck.
% The folder is 98 pages in five batches (1-20, 21-40, 41-60, 61-80, 81-98);
% this is the fourth, done in parallel with the others and aligned afterwards
% on batches 2 and 3 (transcripts/150/batch-02.fr.tex, batch-03.fr.tex).
%
% Pass: Opus 5.5 (claude-opus-5-5), 2026-09-26 — first pass, unchecked against the pages by a human.
%
% Pages skipped: none. Every page carries his mathematics.
%
% Pages summarised: none. Cancelled passages are given as \struck{}: the
% boxed and barred lower half of page 61 (a first « de façon équivalente »
% restatement of the Scholie as an ordered space), the barred head of
% page 62 (a first c)), the boxed start of the conditions on page 68, a
% bracketed first construction of the tubular neighbourhoods on page 73.
% Page 79 is barred entire by two long diagonals; it is transcribed, and the
% cancellation recorded once in a \note (page 80 takes up its (62)-(63)).
%
% His pagination and numbered runs. The batch continues the run headed
% « Espaces stratifiés (étude par voisinages coniques) » (batch 1, sheets
% 1)-5); batch 2, sheets 6-16; batch 3, sheets 17-26). Sheets numbered top
% left: 27 (p. 62), 28 (64), 29 (66), 30 (68), 31 (70), 32 (72), 33 (74),
% 34 (76), 35 (78), 36 (80); the unnumbered odd pages 61-79 are their
% continuations (probably versos). Page 61 continues the « Scholie » of
% batch 3 (sheet 26, p. 60), which broke off at 3° a): it opens with b).
% His cross-references « Scholie (p. 26/27) », « p. 27 », « p. 28 » are
% to his sheets 26-28 (our pp. 60, 62, 64). Page 80 (sheet 36) ends on
% (66) J ≺ L, taken up by batch 5's page 81 (sheet 37, (67)).
% Paragraphs: « 12. Récapitulation » (p. 64, the underlined 12 written at
% the head of the left margin), « 13. Une
% digression » (p. 66; the numeral could be read « B. »), « 14. Description
% heuristique des Sigma_i, Sigma_ij » (p. 72), « 15. Introduction des
% provoisinages » (p. 76); paragraph 16 opens batch 5 (p. 84). Numbered
% formulas (40)-(66), continuing batch 3's (39): (40), (41) the two
% hexagonal diagrams, (42)-(47), a struck (48) and (49) on p. 74 before the
% valid (48)-(49); (49) is used again on p. 75 for Sigma_ij; (50)-(61);
% (62)-(63) on the barred p. 79, and (62)-(66) again on p. 80. Kept as
% written, with \note{}s. Circled letters (a), (b) on p. 70. No date on
% these pages.
%
% What the batch is about. Pages 61-63: the end of the Scholie — the
% cartesian square Sigma*(2) -> Sigma*(1) over the boundary map, the
% transversality of the source map to the bordant immersion, the
% conditions a)-d) on an ordered topological space and the cases
% 1° equisingular, 2° non-singular, 3° both; the hexagonal diagram (40) of
% s, b, s', b', p_2 with its two cartesian squares. Pages 64-65: a
% « récapitulation » of the data of a « déploiement » over a finite ordered
% set I (spaces Sigma_i, Sigma_ij, maps s_ij, b_ij) with conditions a)-e)
% and the hexagon (41) for i<j<k; their translation into Sigma = ⊔ Sigma_i,
% Sigma*(1) = ⊔ Sigma_ij. Pages 66-71: the equisingular and non-singular
% cases; paragraph 13, a digression on a topological category Sigma with s
% fibrant — pi_0(Sigma(0)) is a preordered set (lemma on pi_0 of a
% cartesian square with a fibration), the canonical functor to it, the
% conditions a)-d) which make every arrow inside a component an identity;
% the « conséquence »: equisingular déploiements indexed by I are
% equivalent to an ordered space Sigma plus a strictly increasing map
% pi_0(Sigma) -> I; a « Remarque » comparing the two contexts of
% déploiement met so far. Pages 72-75: paragraph 14, a heuristic
% construction of Sigma_i, Sigma_ij by standard tubular neighbourhoods
% (42)-(43), made coherent by a dimension function delta : I -> N, the
% filtrations (44)-(48) and the formulas (49)-(50) for Sigma_{i_0...i_n}.
% Pages 76-80: paragraph 15, « provoisinages » (pro-neighbourhoods) \mathcal V_i of the strata X_i^*,
% V_ij = V_i ∩ V_j, X as the colimit (53), relative versions V^Z (55)-(57),
% X_J^* for « idéaux » and « bandes » J (58)-(61), and the provoisinages
% V_J^L (62)-(66) with the domination relation J ≺ L.
%
% Notation, following batches 2-3: Sigma^{*}(n) for the spaces of strict
% flags, \underline s, \underline b, \underline s', \underline b', p_1, p_2
% as he underlines them; \underline{\Sigma} for the topological category
% (he doubles the underline); \mathring T for interiors, \dot T for
% boundaries, as in batch 2; \mathcal V (his round V) for provoisinages, as batch 5 does; \widetilde J, \widetilde I for his tildes; \varinjlim
% for his « lim » over an arrow. « provoisinage » is read from the
% interlinear « pro » on pp. 76 and 79 (batch 5 has pro-tubular
% neighbourhoods); (h) over an arrow, his mark for a homotopy equivalence,
% is set \overset{(h)}{...} as in batch 5. Spelling: he writes « côniques » in
% batch 2; no occurrence of the word in this batch. His underlinings in
% running text are \emph{}; circled letters are (a) with a \note.
%
% Hardest pages: 61 (the barred lower half and the sideways margin notes),
% 62 and 64 (long sideways notes in the left margin), 65 (the last lines),
% 79 (barred, and margin notes). The formulas and numbered statements carry;
% most of the sideways marginalia is read only in fragments and left \ill{}.

\documentclass[11pt,a4paper]{article}
% Le préambule commun à toutes les transcriptions du fonds Grothendieck.
%
% Il tient en une idée : **l'incertitude se compose, elle ne se résout pas.**
% Une transcription qui lisse un mot douteux a détruit la seule chose pour
% laquelle on la faisait — un lecteur doit pouvoir savoir, à chaque endroit,
% ce qui était sur le feuillet et ce qui a été ajouté. Les six macros qui
% suivent sont donc l'appareil critique, et non de la décoration.
%
% Les mêmes commandes sont reconnues par `scripts/render.mjs`, qui produit la
% vue de lecture du site. Une macro ajoutée ici doit l'être aussi là-bas,
% faute de quoi le rendu échouera bruyamment — ce qui est voulu : un rendu
% silencieusement faux est pire qu'un rendu absent.

\ProvidesPackage{grothendieck}[2026/08/08 Transcriptions du fonds Grothendieck]

\usepackage{amsmath,amssymb,amsthm}
\usepackage{mathrsfs}
\usepackage{stmaryrd}
% Les diagrammes commutatifs sont partout dans ce fonds : c'est la langue
% même de la géométrie algébrique de Grothendieck, et une transcription qui
% les rendrait en prose ne transcrirait rien.
\usepackage{tikz-cd}
% Français par défaut — c'est la langue des feuillets — mais l'anglais est
% chargé aussi : la traduction et le résumé sont des documents anglais, et la
% typographie française (espaces avant les deux-points) y serait une faute.
% Ces documents basculent par \selectlanguage{english} en tête de corps.
\usepackage[english,french]{babel}
\usepackage{fontspec}
\usepackage{geometry}
\geometry{a4paper,margin=25mm,marginparwidth=28mm}
\usepackage{marginnote}
\usepackage{xcolor}
% ulem, not soul. `soul`'s \st and \ul are built on a character-by-character
% scan that XeLaTeX's font machinery breaks: the document fails with an
% undefined control sequence rather than degrading. ulem does the same two jobs
% and is Unicode-engine safe. `normalem` stops it hijacking \emph, which the
% transcriptions use for Grothendieck's own emphasis.
\usepackage[normalem]{ulem}
\usepackage{hyperref}
\hypersetup{colorlinks=true,linkcolor=black,urlcolor=black,pdfborder={0 0 0}}

% --- Métadonnées du lot -----------------------------------------------------
% Reprises telles quelles par le script de rendu, qui les affiche en tête de
% la vue de lecture. Elles fixent ce que le fichier prétend couvrir : sans
% elles, une transcription flotte sans point d'attache dans un fonds de seize
% mille feuillets.

\newcommand{\folder}[1]{\gdef\grothFolder{#1}}
\newcommand{\batch}[1]{\gdef\grothBatch{#1}}
\newcommand{\leaves}[2]{\gdef\grothFirst{#1}\gdef\grothLast{#2}}
\newcommand{\foldertitle}[1]{\gdef\grothFoldertitle{#1}}
\gdef\grothFolder{?}\gdef\grothBatch{?}\gdef\grothFirst{?}\gdef\grothLast{?}\gdef\grothFoldertitle{}

% --- Le résumé --------------------------------------------------------------
%
% Il ouvre la lecture modernisée, sous le titre « Résumé », et il est composé
% en petit. Ce n'est pas un ornement : c'est le seul passage du document qui
% ne soit pas des mathématiques, et un lecteur doit pouvoir le reconnaître
% avant de l'avoir lu — puis le sauter s'il connaît déjà le sujet, ou s'y
% arrêter s'il ne le connaît pas. Le filet qui le suit marque la reprise du
% corps.

\newenvironment{resume}
  {\small\setlength{\parskip}{0.45em}}
  {\par\medskip\hrule\medskip}

% --- Mots-clés ---------------------------------------------------------------
%
% En anglais, en clôture du résumé de la lecture modernisée : le vocabulaire
% moderne sous lequel chercher ce que les feuillets construisent. Le manifeste
% du site (`npm run manifest`) les extrait pour étiqueter la cote entière —
% cette ligne est la source unique des tags, il n'y a pas de fichier à côté.
\newcommand{\keywords}[1]{\par\smallskip\noindent
  {\footnotesize\textsc{Keywords} --- \textit{#1}}\par}

% --- L'appareil critique ----------------------------------------------------

% Le numéro de feuillet, en marge. C'est l'ancre qui permet de régler un
% désaccord sur une formule en regardant *un* feuillet plutôt qu'une chemise
% de six cents. Le site s'en sert aussi pour tourner les pages du fac-similé
% quand on fait défiler la transcription.
\newcommand{\leaf}[1]{%
  \marginnote{\footnotesize\color{gray!70}\textsf{#1}}%
  \label{leaf:#1}\ignorespaces}

% Illisible : on ne devine pas. Un mot inventé qui se lit comme les autres est
% le pire résultat possible.
\newcommand{\ill}{\textcolor{red!60!black}{[\,\ldots\,]}}

% Lecture douteuse : le mot est donné, et signalé comme incertain.
\newcommand{\uncertain}[1]{\uline{#1}}

% Ajout de l'éditeur — une lettre manquante, un mot sous-entendu.
\newcommand{\add}[1]{\textcolor{blue!50!black}{[#1]}}

% Biffé par Grothendieck lui-même. Ce qu'il a rayé fait partie du document :
% une rature dit souvent où la pensée a changé de direction.
\newcommand{\struck}[1]{\sout{#1}}

% Note du transcripteur, sur l'état matériel du feuillet.
\newcommand{\note}[1]{\par\smallskip{\small\color{gray!60!black}\textit{[#1]}}\par\smallskip}

% Note marginale de Grothendieck, distincte du corps du texte.
\newcommand{\marginal}[1]{\par\smallskip\noindent\hspace{1em}%
  {\small\color{gray!60!black}#1}\par\smallskip}

% --- Titre ------------------------------------------------------------------

\AtBeginDocument{%
  \begin{center}
    {\large\bfseries Fonds Alexandre Grothendieck --- cote n\textsuperscript{o}~\grothFolder}\\[2pt]
    {\itshape\grothFoldertitle}\\[4pt]
    {\small feuillets \grothFirst--\grothLast\ (lot \grothBatch)}\\[2pt]
    {\footnotesize\color{gray!60!black}Transcription établie d'après le fac-similé numérisé
     par l'Université de Montpellier.\\
     Les lectures incertaines sont soulignées, les illisibles notées [\,\ldots\,],
     les ajouts entre crochets.}
  \end{center}
  \vspace{1em}}

\endinput


\folder{150}
\batch{4}
\pages{61}{80}
\dating{1981-1982}
\watermark{Édition de démonstration}
\foldertitle{Espaces stratifiés et voisinages côniques (ou : déploiement des espaces stratifiés) : notes manuscrites (s.d.)}

\begin{document}

\page{61}
\note{Sans numéro de sa main : suite du feuillet 26 (page 60). La page poursuit la « Scholie » de la page 60, interrompue à 3°) a) : elle commence par b).}

b) L'application
\[
\Sigma^{*}(2)\xrightarrow{\ (p_1\circ\underline s',\ p_1\circ\underline b')\ }\Sigma\times\Sigma
\]
se factorise par
\[
\Sigma^{*}(1)\xrightarrow{\ (p_1,\underline b)\ }\Sigma\times\Sigma
\]
en un morphisme
\[
p_2 : \Sigma^{*}(2)\longrightarrow\Sigma^{*}(1)
\]
qui rend le carré

\begin{tikzcd}
\Sigma^{*}(2) \arrow[r, hook, "\underline b'"] \arrow[d, "p_2"'] & \Sigma^{*}(1) \arrow[d, "p_1"] \\
\Sigma^{*}(1) \arrow[r, hook, "\underline b"] & \Sigma
\end{tikzcd}
\note{Sur la page, $p_2$ descend en diagonale sous $\underline b'$ ; les deux flèches horizontales ont une queue d'inclusion. L'étiquette de la flèche du haut, $(p_1\circ\underline s',\ p_1\circ\underline b')$, est surchargée : la lecture du second terme est incertaine.}

\emph{cartésien}, et $p_1 : \Sigma^{*}(1)\to\Sigma$ est transversal à l'\struck{\ill{}} immersion-bord $\underline b : \Sigma^{*}(1)\to\Sigma$.

\marginal{\emph{NB} Si $\Sigma\to$ \ill{} \ldots\ $\Sigma\times\Sigma$ \ldots\ $\Sigma^{*}(1)\xrightarrow{\underline b}\Sigma$ \ldots\ il \ill{} \ldots\ $\Sigma^{*}(1)\to\Sigma\times\Sigma$ \ldots\ \ill{}}\note{notes en travers dans la marge gauche, en face de b) ; une seconde, plus bas, commence par « cette condition \ldots\ transversalité \ldots\ » ; lecture très incertaine.}

\struck{De façon équivalente, une telle structure est donnée par un espace topologique \add{ordonné} $\Sigma$, muni d'une relation d'ordre $x\le y$ satisfaisant les conditions suivantes \struck{dont le graphe est} a) le graphe \struck{$\Sigma^{*}(1)$} $\Sigma(1)$ de cette relation est fermé dans $\Sigma\times\Sigma$, et $\Sigma(1)\xrightarrow{\underline b}\Sigma$ \ill{} immersion \ill{}. b) Si $\Sigma^{*}(1)$ est le graphe de la relation d'ordre \add{stricte}, $\Sigma^{*}(1)\xrightarrow{\underline b}\Sigma$ est une immersion \add{bordante} propre (donc les \ill{} $(=\Sigma^{*}(1)_x)$). a) Si $x\in\Sigma$, $\Sigma_{<x}$ est un ens.\ (fini) tot.\ ordonné.}\note{Toute la moitié inférieure de la page, à partir de « De façon équivalente », est encadrée et barrée de longues diagonales ; lecture incertaine dans le détail. Elle est reprise, sous une autre forme, en tête de la page 62.}

\page{62}
\note{En haut à gauche, « 27 » de sa main.}

\struck{c) Si $(z_\alpha)_{\alpha\in I}$ est l'une des éléments de $\Sigma_{<x}=\underline b^{-1}(\lbrace x\rbrace)$ (où $\underline b : \Sigma^{*}(1)\to\Sigma$), les \ill{}}\note{Ces deux lignes sont barrées de deux diagonales et encadrées.}

a) Pour tout $x\in\Sigma$, l'ens.\ $\Sigma_{<x}$ \struck{(fini, \ldots} \struck{$\simeq\underline b^{-1}(\lbrace x\rbrace)$} \add{tot\textsuperscript{t}} ordonné

b) Si $\Sigma^{*}(1)$ est le graphe de la relation $x<y$, et $\underline s,\underline b : \Sigma^{*}(1)\rightrightarrows\Sigma$ sont les applications source et but, l'application $\underline b$ est une immersion \struck{\ill{}} \uncertain{propre} (donc les $\Sigma_{<x}=\underline b^{-1}(\lbrace x\rbrace)$ sont finis)

c) Si $(z_\alpha)_{\alpha\in I}$ sont les éléments de $\Sigma_{<x}=\underline b^{-1}(\lbrace x\rbrace)$, alors \add{la famille} des \struck{\ill{}} germes \add{en $x$} de sous-espaces \struck{bordants} \add{de $\Sigma$}, images par $\underline b$ des $(\Sigma^{*}(1),z_\alpha)$, est \add{une famille} \struck{\ill{}} \add{transversale} de sous-espaces \emph{bordants}.

d) $\underline s : \Sigma^{*}(1)\to\Sigma$ est \emph{transversal} p.r.\ à l'immersion bordante $\underline b : \Sigma^{*}(1)\to\Sigma$. \struck{\ill{}}\note{La lettre devant « : » est grasse et surchargée ($\underline s$ ou $p_1$) ; le but est écrit $\Sigma^{*}$, avec une étoile qu'on ne sait pas lire autrement.}

\marginal{au sens \ill{} d'une \ill{} \ldots\ pour les $z_\alpha\in\Sigma^{*}(1)$ \ldots\ $\in\Sigma^{*}$ \ldots}\note{note en travers dans la marge gauche, en face de d) ; lecture très incertaine.}

1°) Cas équisingulier : $\underline s : \Sigma^{*}(1)\to\Sigma$ fibrant \add{(propre)}

2°) Cas non singulier : $\Sigma$ est une variété à bord, dont le bord est $\underline b(\Sigma^{*}(1))$.

3°) Cas équisingulier-non singulier : $\underline s : \Sigma^{*}(1)\to\Sigma$ est \add{propre} fibrant, à fibres des variétés à bord, et le fibré des bords est l'image de $\Sigma^{*}(2)$ par le morphisme $\Sigma^{*}(2)\xrightarrow{\underline s'}\Sigma^{*}(1)$ \struck{\ill{}} \add{donné} par $(x<y<z)\longmapsto(x<z)$.

\marginal{c'est sans doute conséquence de la conjonction des deux conditions 1°), 2°)}\note{en travers dans la marge gauche, en face de 3°).}

\page{63}
\note{Sans numéro de sa main : suite du feuillet 27.}

Les morphismes les plus importants dans la situation apparaissent dans le diagramme \add{hexagonal suivant}

(40)

\begin{tikzcd}[column sep=small]
 & \Sigma^{*}(1) \arrow[r, hook, "\underline b"] \arrow[dl, "\underline s"'] & \Sigma & \\
\Sigma & \Sigma^{*}(2) \arrow[u, hook, "\underline s'"] \arrow[rr, hook, "\underline b'"] \arrow[d, "p_2"'] & & \Sigma^{*}(1) \arrow[ul, hook, "\underline b"'] \arrow[dl, "\underline s"] \\
 & \Sigma^{*}(1) \arrow[r, hook, "\underline b"] \arrow[ul, "\underline s"] & \Sigma &
\end{tikzcd}
\note{Hexagone de sa main, le numéro (40) à gauche ; les deux carrés, en haut et en bas de $\Sigma^{*}(2)$, portent chacun la mention « cart ». Le $\Sigma^{*}(1)$ de droite est au milieu de la hauteur, entre les deux $\Sigma$ du côté droit.}

où les flèches marquées $\hookrightarrow$ (en fait $\underline b$, \struck{\ill{}}, $\underline b'$, $\underline s'$) sont des immersions bordantes \add{(propres)}, et les flèches marquées $\longrightarrow$ (en fait, $\underline s$ et $p_2$) sont (dans le cas équisingulier) fibrantes et propres (voir, de plus : fibres variétés à bord, dont le bord est dans le cas de $\underline s : \Sigma^{*}(1)\to\Sigma$ l'image de $\Sigma^{*}(2)\xrightarrow{\underline s'}\Sigma^{*}(1)$ \add{i.e.\ l'ens.\ des $(x<y)$ tels que \uncertain{$y$ ne soit pas successeur de $x$}}). On fera attention que les applications fibrantes $\underline s$, $p_2$ n'ont \struck{aucune} \add{rien} de surjectif -- l'image de $\underline s$ est formée des $x\in\Sigma$ qui ne sont pas \emph{maximaux}, celle de $p_2$ est formée des $(x<y)$ tels que $y$ ne soit pas maximal.

\marginal{dans le cas de $p_2 : \Sigma^{*}(2)\to\Sigma^{*}(1)$, le fibré des bords est l'image de $\Sigma^{*}(3)$ \ldots}\note{en travers dans la marge gauche ; lecture incertaine.}

\page{64}
\note{En haut à gauche, « 28 » de sa main.}

\section*{Récapitulation}

12.\note{Numéro souligné, écrit au-dessus et à gauche du titre, au début de la note en travers de la marge.} \emph{Récapitulation} (des données de déploiement, sur un ens.\ ordonné $I$, fini disons\ldots)

a) Espaces $\Sigma_i$ ($i\in I$)

b) Espaces $\Sigma_{ij}$ ($i<j$, $i,j\in I$)

c) Applications $\Sigma_{ij}\xrightarrow{s_{ij}}\Sigma_i$, $\Sigma_{ij}\xrightarrow{b_{ij}}\Sigma_j$

\marginal{C'est la structure \ill{} envisagée dans les Scholie (p.~26/27) plus une application strictement croissante $\Sigma\to I$ \ldots\ injection \ldots\ $\Sigma_{<x}$ \ldots\ $I_{<i}$ \ldots}\note{longue note en travers dans la marge gauche, en face du titre et de a)-c) ; « p. 26/27 » renvoie à ses feuillets 26 et 27 (pages 60 et 62) ; lecture incertaine dans le détail.}

\emph{Conditions}

a) $\Sigma_{ij}\longrightarrow\Sigma_j$ \struck{\ill{}} \uncertain{plongement} fermé \struck{bordant}

\marginal{[On identifie $\Sigma_{ij}$ à un sous-espace fermé de $\Sigma_j$ par $b_{ij}$]}\note{entre crochets, à droite de a).}

b) $\Sigma_{ik}\cap\Sigma_{jk}=\emptyset$ si $i,j$ non comparables \struck{\ill{}} (\emph{NB} $i,j<k$)

c) Pour $i<j<k$ posons $\Sigma_{ijk}=\Sigma_{ik}\cap\Sigma_{jk}\subset\Sigma_k$. On a alors dans $\Sigma_{jk}$ :
\[
\Sigma_{ijk}=b_{ij}^{-1}(\Sigma_{ij})
\]
\note{Égalité soulignée et encadrée ; un premier membre et un « $\Sigma_{ik}\cap$ » sont biffés. L'indice de $b^{-1}$ est lu tel quel ; on attendrait $s_{jk}^{-1}$.}

\marginal{\emph{NB} La condition b) équivaut à dire que dans l'ens.\ ordonné $\Sigma=\coprod\Sigma_i$, les $\Sigma_{<x}$ sont tot\textsuperscript{t} ordonnés}\note{en travers dans la marge gauche, en face de b) et c).}

On a donc, pour $i<j<k$, un diagramme commutatif hexagonal

(41)

\begin{tikzcd}[column sep=small]
 & \Sigma_{ik} \arrow[r, hook, "b_{ik}"] \arrow[dl, "s_{ik}"'] & \Sigma_k & \\
\Sigma_i & \Sigma_{ijk} \arrow[u, hook] \arrow[rr, hook] \arrow[d] & & \Sigma_{jk} \arrow[ul, hook, "b_{jk}"'] \arrow[dl, "s_{jk}"] \\
 & \Sigma_{ij} \arrow[r, hook, "b_{ij}"] \arrow[ul, "s_{ij}"] & \Sigma_j &
\end{tikzcd}
\note{Même disposition que (40), page 63 ; les deux carrés portent « cart ».}

où les flèches $\hookrightarrow$ sont des \add{\uncertain{cf.\ plus bas}} immersions fermées (bordantes), et les flèches $\longrightarrow$ sont des applications qui seront fibrantes et propres dans le cas équisingulier.

\page{65}
\note{Sans numéro de sa main : suite du feuillet 28.}

d) Les $\Sigma_{j}\hookrightarrow\Sigma_{ji}$ sont bordants, et si $x\in\Sigma_i$, \struck{les} dans les $\Sigma_{ji}\ni x$ (\struck{pour $j$} \add{indexés} $j<i$ \struck{tels que l'on $\Sigma$}. $I_{<i,x}\overset{\text{déf}}{=}\lbrace j\in I\mid j<i,\ \Sigma_{ji}\ni x\rbrace$, qui est totalement ordonné par b)) sont des sous-espaces-bords transversaux en $x$ -- i.e.\ les $\Sigma_{ji}$ ($j\in I_{<i}$) sont des sous-espaces bords de $\Sigma_i$ mutuellement transversaux.\note{« $\Sigma_{j}\hookrightarrow\Sigma_{ji}$ » tel quel ; on attend $\Sigma_{ji}\hookrightarrow\Sigma_i$.}

e) $\Sigma_{ij}\xrightarrow{s_{ij}}\Sigma_i$ est transversal p.r.\ au système transversal de morphismes \struck{bordants} $b_{ki} : \Sigma_{ki}\to\Sigma_i$.

Si on pose
\[
\Sigma=\coprod_{i\in I}\Sigma_i,\qquad \Sigma^{*}(1)=\coprod_{\substack{i,j\in I\\ i<j}}\Sigma_{ij}
\]
et $\underline s,\underline b : \Sigma^{*}(1)\to\Sigma$ induits par les $s_{ij}$, $b_{ij}$, alors les conditions a) b) c) signifient que $(\underline s,\underline b) : \Sigma^{*}(1)\to\Sigma\times\Sigma$ est inj.\ et a comme \struck{graphe} \add{image} le graphe d'une relation d'ordre strict, $\underline b : \Sigma^{*}(1)\to\Sigma$ \add{immersion} propre, et que l'application \add{\ill{}} $\varphi : \Sigma\to I$ est strictement croissante, et que pour tout $x\in\Sigma$, l'application induite $\varphi : \Sigma_{<x}\to I_{<i}$ ($i=\varphi(x)$) est \add{injective}, \ill{} que $\Sigma_{<x}$ est tot.\ ordonné.\note{Dans « $\varphi : \Sigma\to I$ », $\Sigma$ porte un exposant surchargé ; le $I$ du but est gras. La fin de la page est serrée, les ajouts interlinéaires se lisent mal.}

\page{66}
\note{En haut à gauche, « 29 » de sa main.}

Enfin, la condition d) signifie que cette structure \add{d'espace top.\ ordonné} (satisfait la condition c) du Scholie (p.~27)).\note{« p. 27 » : son feuillet 27, page 62.}

\emph{Cas équisingulier} : les $s_{ij}$ sont fibrants propres

\emph{Cas non singulier} : les $\Sigma_i$ sont des variétés topologiques, le bord de $\Sigma_i$ est la réunion des $\Sigma_{ji}$ pour $j<i$ \struck{(NB il en suffit} \struck{\ill{}}

\marginal{On voudra que les $\Sigma_{ij}$ soient des variétés à bord, par la réunion des \ill{} $\Sigma_{kij}$ \ldots\ mais \ill{} des conditions \ill{} \ldots}\note{en travers dans la marge gauche ; lecture incertaine.}

\emph{Cas équisingulier-non singulier} $s_{ij} : \Sigma_{ij}\to\Sigma_i$ est une fibration propre, à fibres des variétés à bord, le fibré des bords est réunion des $\Sigma_{ijk}$ pour $i<j<k$.\note{L'indice $ij$ de la source est surchargé, et le « $j$ » de « $i<j<k$ » est biffé puis récrit.}

\section*{Une digression}

13.\note{Numéro lu « 13. » (le 1 et le 3 accolés pourraient se lire « B. ») ; le précédent est 12. (page 64), les suivants 14. (page 72) et 15. (page 76).} \emph{Une digression}. Revenons au cas de $\Sigma$, $\Sigma^{*}(1)\subset\Sigma\times\Sigma$, sans données de $I$. Plus généralement, considérons une \emph{catégorie topologique}
\[
\begin{aligned}
\underline{\Sigma}=\Bigl(&\Sigma(0)=\mathrm{Ob}(\underline{\Sigma}),\ \Sigma(1)=\mathrm{Fl}(\underline{\Sigma}),\ \Sigma(1)\overset{\underline s}{\underset{\underline b}{\rightrightarrows}}\Sigma(0),\\
&(\Sigma(1),\underline b)\times_{\Sigma(0)}(\Sigma(1),\underline s)\overset{\text{déf}}{\Longrightarrow}\Sigma(2)\xrightarrow[\text{compos.\ des flèches}]{\text{compl.}}\Sigma(1)\Bigr),
\end{aligned}
\]
et supposons que $\underline s : \Sigma(1)\to\Sigma(0)$ soit \emph{fibrant}. Je dis qu'alors $\pi_0(\Sigma(0))=I(\underline{\Sigma})$ est \emph{un ens.\ \add{pré}ordonné}, en prenant comme graphe de la relation d'ordre l'image de
\[
\pi_0(\Sigma(1))\xrightarrow{\ \pi_0(\underline s,\underline b)\ }\pi_0(\Sigma(0)\times\Sigma(0))\simeq\pi_0(\Sigma(0))\times\pi_0(\Sigma(0))
\]
\note{$\underline{\Sigma}$ est souligné deux fois. L'étiquette au-dessus de la flèche vers $\Sigma(1)$ se lit « compl. » ou « compo. », sous la flèche « composition des flèches ». Le « pré » d'« ens. préordonné » est écrit au-dessus d'un mot biffé. Le $\simeq$ final est écrit vertical, sous $\pi_0(\Sigma(0)\times\Sigma(0))$.}

\page{67}
\note{Sans numéro de sa main : suite du feuillet 29.}

\struck{Cela posé} Cette relation est évidemment \struck{symétrique} \add{réflexive}, et elle est transitive grâce au fait que
\[
\pi_0(\Sigma(2))\longrightarrow\bigl(\pi_0(\Sigma(1)),\pi_0(\underline b)\bigr)\times_{\pi_0(\Sigma(0))}\bigl(\pi_0(\Sigma(1)),\pi_0(\underline s)\bigr)
\]
est surjective, ce qui est facile à prouver, grâce à l'hypothèse que $\underline b$ est fibrant, par le\note{« $\underline b$ » tel quel ; l'hypothèse posée page 66 porte sur $\underline s$.}

\emph{Lemme} : Soit

\begin{tikzcd}
X \arrow[d, "f"'] & X' \arrow[l, "f'"'] \arrow[d, "g'"] \\
Y & Y' \arrow[l, "g"]
\end{tikzcd}

un carré cartésien d'espaces top., avec $f$ fibrant, alors $\pi_0(X')\to\pi_0(X)\times_{\pi_0(Y)}\pi_0(Y')$ est surjectif.

Ceci posé, considérons la catégorie top.\ discrète associée à l'ens.\ \add{pré}ordonné $I$ \add{$=\pi_0(\Sigma_0)$}, que par abus de langage \add{\ill{}} on désigne par $I$. On a alors un foncteur \add{canonique} de catégories topologiques \struck{canoniques}
\[
\underline{\Sigma}\xrightarrow{\ f\ }\underline I
\]
qui est donc donné par des applications continues
\[
\left\lbrace
\begin{array}{l}
\Sigma(0)\xrightarrow{\ \varphi_0\ }I_0=I\\
\Sigma(1)\xrightarrow{\ \varphi_1\ }I_1=\text{graphe de la relation d'ordre}\\
\qquad=\text{image de }\pi_0(\Sigma_1)\text{ dans }\pi_0(\Sigma_0)\times\pi_0(\Sigma_0)=I\times I
\end{array}\right.
\]
\note{Sous « image de $\pi_0(\Sigma_1)$ », une première fin « dans $\pi_0(\Sigma_0)\times\Sigma_1$ » est biffée.}

i.e.\ des décompositions en sommes directes
\[
\Sigma(0)=\coprod_{i\in I}\Sigma_i,\qquad \Sigma(1)=\coprod_{\substack{i,j\in I\\ i\le j}}\Sigma_{ij}
\]
compatibles avec $\underline s$, $\underline b$, i.e.\ $\underline s : \Sigma_{ij}\to\Sigma_i$, $\underline b : \Sigma_{ij}\to\Sigma_j$, avec
\[
\Sigma(0)\xrightarrow{\ \delta\ }\Sigma(1)\qquad(\delta(x)=\mathrm{id}_x^{\underline{\Sigma}})
\]
appliquant $\Sigma_i$ dans $\Sigma_{ii}$, \struck{\ill{}} (ce qui indique déjà la compatibilité de $\varphi=(\varphi_0,\varphi_1)$ avec la composition). De plus, $\varphi$ \uncertain{a} une

\marginal{\emph{NB} \ldots\ \ill{} \ldots\ $\Sigma_0$ et $\Sigma_1$ \ldots}\note{en travers dans la marge gauche ; illisible pour l'essentiel.}

\page{68}
\note{En haut à gauche, « 30 » de sa main.}

\emph{propriété universelle} évidente pour les morphismes de \struck{la} catégories top.\ $\underline{\Sigma}$ dans un ens.\ \add{pré}ordonné.

§ Supposons maintenant que $\underline{\Sigma}$ satisfait les conditions \struck{: toute flèche inversible de $\underline{\Sigma}$ est une identité,} \struck{a) l'application $\delta : \Sigma_0\to\Sigma_1$ est un plongement ouvert et fermé. Ceci est le cas si p.ex.\ $\underline b : \Sigma_1\to\Sigma_0$ est une immersion propre \ldots)}\note{Le début des conditions est encadré et barré de diagonales ; le « § » est de sa main.}

a) $\Sigma(1)\xrightarrow{\underline b}\Sigma(0)$ est une immersion propre ($\Rightarrow\delta : \Sigma(0)\to\Sigma(1)$ est une immersion ouverte et fermée, de sorte que \struck{$\Sigma(1)=\delta(\Sigma(0))\amalg$}
\[
\Sigma(1)\simeq\delta(\Sigma(0))\amalg\Sigma^{*}(1)\ )
\]

b) $\Sigma^{*}(1)\xrightarrow{\underline b|\Sigma^{*}(1)}\Sigma(0)$ est une immersion à images localement \emph{rares} i.e.\ : intérieurs vides\ldots

c) Toute composante connexe de $\Sigma(0)$ est de dim.\ finie.

Alors on trouve que
\[
\Sigma^{*}_{ii}\overset{\text{déf}}{=}\Sigma_{ii}\cap\Sigma^{*}(1)=\emptyset ,
\]
pour des raisons de dimension (car $\Sigma^{*}_{ii}$ est fibré sur $\Sigma_i$, donc s'il est non vide, de $\dim\ge\dim\Sigma_i$, mais comme $\Sigma^{*}_{ii}\to\Sigma_i$ est une immers.\ rare, on doit avoir aussi $\dim\Sigma^{*}_{ii}<\dim\Sigma_i$, donc $\Sigma^{*}_{ii}$ est vide), i.e.
\[
\boxed{\Sigma_i\simeq\Sigma_{ii}}
\]
: toute flèche \struck{dans} entre deux \struck{objets} éléments de $\Sigma(0)$ qui sont dans une même composante connexe \struck{est} est une identité.

\page{69}
\note{Sans numéro de sa main : suite du feuillet 30.}

Supposons en plus

d) toute flèche \emph{inversible} de $\underline{\Sigma}$ est une identité (p.ex.\ $\underline{\Sigma}$ est une catégorie définie par un ens.\ \emph{ordonné}).

Alors on conclut aisément que l'ens.\ préordonné $I$ est \emph{ordonné}, et que toute flèche de $\underline{\Sigma}$ qui n'est pas une identité i.e.\ $u\in\Sigma^{*}(1)$, a son image dans $I(1)$ qui n'est pas une identité ; \struck{i.e.} si $\underline{\Sigma}$ est associé à un \add{espace top.} \struck{ens.\ ord.} ordonné \add{$\Sigma=\Sigma(0)$}, cela signifie que $\Sigma\to I$ est \emph{strictement} croissante.

Revenons alors à un \struck{catég} ens.\ ordonné topologique
\[
\underline{\Sigma}=\bigl(\Sigma(0),\ \Sigma(1)=\delta(\Sigma(0))\amalg\Sigma^{*}(1)\bigr),\qquad \Sigma(1)\subset\Sigma\times\Sigma ,
\]
satisfaisant aux conditions a) b) c) et 1°) \add{au moins} de la page (27) (correspondant au déploiement d'une stratification localement équisingulière), et tel que les comp.\ connexes de $\Sigma$ \struck{soient} \add{\ill{}} de dim.\ finie. Considérons alors la \struck{factorisation} \add{\ill{}} canonique
\[
\underline{\Sigma}\longrightarrow\underline I\qquad\text{avec } I=\pi_0(\Sigma),
\]
et supposons de plus que pour $i<j$ ($i,j\in I$)
\[
\Sigma_{ij}\longrightarrow\Sigma_j\quad\text{soit \emph{mono} i.e.\ plongement fermé.}
\]
On est alors sous les conditions de p.~28, données a) b) c) $\bigl(\Sigma_i,\ \Sigma_{ij},\ s_{ij},b_{ij} : \Sigma_{ij}\rightrightarrows\Sigma_i,\Sigma_j\bigr)$\note{« page (27) », « p. 28 » : ses feuillets 27 et 28 (pages 62 et 64). Sous $\Sigma(1)$, un « $\cap$ » vertical renvoie à $\Sigma\times\Sigma$.}

\page{70}
\note{En haut à gauche, « 31 » de sa main.}

satisfaisant les conditions a) c), d) de p.~28.

Montrons aussi la condition b) de p.~28
\[
\Sigma_{ik}\cap\Sigma_{jk}=\emptyset\quad\text{si } i,j<k,\ i,j\ \text{non comparables.}
\]
Si on avait $x\in\Sigma_{ik}\cap\Sigma_{jk}\subset\Sigma_k$, donc on aurait deux éléments $x_i$ et $x_j$ de $\Sigma$, au-dessus de $i$ et $j$ respectivement, qui sont $<x$ (au-dessus de $k$), donc (comme $\Sigma_{<x}$ est tot\textsuperscript{t} ordonné) \struck{on aurait} $x$ et $y$ seraient comparables, donc aussi $i$ et $j$, absurde.

\emph{Conséquence} Les structures « déploiements équisinguliers de stratifications, indexées par $I$ », décrites dans la p.~28, équivalent à la structure suivante :

(a)\note{lettre cerclée ; de même (b) plus bas.} Un \struck{\ill{}} espace top.\ ordonné $\Sigma$, satisfaisant \struck{aux} conditions a) b) c) et 1°) de la page 27, et \emph{de plus} : la condition \struck{\ill{}} que, \struck{pour tout} posant $\widetilde I=\pi_0(\Sigma)$, on ait, pour $i,j\in\widetilde I$, $i<j$, que le morphisme \add{bordant} $\Sigma_{ij}\xrightarrow{b_{ij}}\Sigma_j$ soit une immersion propre \emph{mono}

(b) Une application strictement croissante (pas nécessairement surjective)
\[
\widetilde I\longrightarrow I
\]

\marginal{Si les \ill{} \ill{} des $\Sigma_i$ \ill{} de dim.\ finie\ldots}\note{en travers dans la marge gauche, en face de (a) ; lecture incertaine.}

\page{71}
\note{Sans numéro de sa main : suite du feuillet 31.}

\emph{Remarque} Dans deux contextes différents${}^{*}$ de « déploiement » de \struck{\ill{}} stratification, on a rencontré des espaces topologiques ordonnés $\Sigma$, satisfaisant les conditions suivantes

\marginal{${}^{*}$ en deux sens assez différents}\note{en travers dans la marge gauche, rattachée par l'astérisque ; lecture incertaine.}

a) $\Sigma(1)\xrightarrow{\underline s}\Sigma$ \struck{\ill{}} $=\Sigma(0)$ est une fibration propre

b) $\Sigma(1)\xrightarrow{\underline b}\Sigma(0)$ est une immersion propre

c) Si $x,y\in\Sigma$ \struck{sont} dans une même composante connexe, alors $x\le y\Rightarrow x=y$ (i.e.\ si $x,y\in\Sigma$, $x<y$, alors $x,y$ sont dans des composantes connexes différentes \struck{\ill{}})

\marginal{conséquence de a) b) si on suppose que $\Sigma^{*}(1)\xrightarrow{\underline b}\Sigma$ est \ill{} et que les comp.\ connexes de $\Sigma$ \ill{} de dim.\ finie\ldots}\note{en travers dans la marge gauche, en face de c) ; lecture incertaine.}

Dans un contexte de déploiement d'une stratification de type $I$ disons, par $\Sigma=\coprod X_i$,
\[
\Sigma_1=\coprod_{i,j\ \text{tels que}\ X_i\subset X_j}X_i\ ),
\]
$\underline s$ est étale, bien entendu, et \uncertain{rien de plus spécial} pour $\underline b$. Dans le cas de la théorie \struck{présente} \add{développée} ici des déploiements \add{de stratifications équisingulières}, $\underline b$ est bordant, et rien de spécial sur $\underline s$. Dans l'un et l'autre cas, on peut appliquer la sortie des $\widetilde I=\pi_0(\Sigma)$ -- dans le premier cas, cela revient à remplacer la stratification par les $X_i$ par celle par les comp.\ connexes des $X_i$ (donnant un $\widetilde I$) ; dans le deuxième, \struck{\ill{}} itou pour les
\[
X_i^{*}=X_i\setminus\bigcup_{j<i}X_j .
\]

\section*{Description heuristique des $\Sigma_i$, $\Sigma_{ij}$}

\page{72}
\note{En haut à gauche, « 32 » de sa main.}

14. \emph{Description heuristique des $\Sigma_i$, $\Sigma_{ij}$}. (Peut se préciser considérablement dans un contexte de topologie modérée). Pour construire $\Sigma_i$, on prend un voisinage tubulaire \add{« standard » $T_{<i}$} \struck{\ill{}} de $\dot X_i=\bigcup_{k<i}X_k$, et on prend
\[
(42)\qquad \Sigma_i=X_i\setminus X_i\cap\mathring T_{<i}\qquad \mathring T_{<i}\overset{\text{déf}}{=}T_{<i}\setminus\dot T_{<i} .
\]
(\emph{NB} Il suffirait en fait ici de prendre le voisinage tubulaire standard dans $X_i$ lui-même, et de prendre donc $\Sigma_i=X_i\setminus\mathring T_{<i}$.)\note{Le point sur $\dot X_i$ surmonte une lettre surchargée. « (42) » est écrit en marge.}

\marginal{\emph{NB} On prendra $T_{<i}=\bigcup_{k<i}T_k$, où $T_k$ vois.\ tub.\ standard \ill{}}\note{en travers dans la marge gauche, en face de (42) ; lecture incertaine.}

Pour construire $\Sigma_{ij}$, on prend un voisinage tubulaire $T'_{ij}$ de $\Sigma_i$ dans $X_j$ \ill{} $\mathring T_{<i}=X'_i$ (\emph{NB} cette fois-ci il faut \struck{prendre} d'abord prendre $T_{<i}$ vois.\ tub.\ standard \emph{dans} $X$, ou sinon, dans $X_j$). On prend $T'_{ij}$ « suffisamment petit » p.r.\ au choix de $T_{<i}$. Puis on prend $\dot T'_{ij}$, et dans $\dot T'_{ij}$ la trace des $\bigcup_{i<k<j}X_k$ (donc c'est vide si $j$ est un successeur de $i$), soit $S\dot T'_{ij}$ (« partie singulière » de $\dot T'_{ij}$). On prend donc un voisinage tubulaire standard \add{$T(S\dot T'_{ij})$} \struck{de} de $S\dot T'_{ij}$ dans $\dot T'_{ij}$, qu'on peut supposer induit par la réunion \add{$\bigcup T_k$} des voisinages tubulaires standard des $X_k$ ($i<k<j$) dans $X$ \ill{} $X_j$ (pris « assez petits » p.r.\ aux choix déjà faits de $T_{<i}$ et $T'_{ij}$). On \struck{prendra} \add{aura} alors
\[
(43)\qquad \Sigma_{ij}=\dot T'_{ij}\setminus\mathring T(S\dot T'_{ij})
\]
\note{Dans « dans $X_j$ \ldots\ $\mathring T_{<i}=X'_i$ », une lettre surchargée précède $\mathring T_{<i}$ ; l'indice $i$ de $X'_i$ est lu tel quel. En marge gauche de la seconde moitié de la page, un croquis : un disque avec son centre marqué, d'où partent trois traits horizontaux parallèles.}

\page{73}
\note{Sans numéro de sa main : suite du feuillet 32.}

de sorte qu'on ait
\[
\Sigma_{ij}\subset\Sigma_j
\]
à condition d'avoir ajusté de façon évidente les choix qui donnent $\Sigma_i$, $\Sigma_j$.

\struck{Cette construction donne}

Pour ajuster ces choix, de façon à avoir, \struck{étant} $i<j<k$, des $\Sigma_{ik}$ et $\Sigma_{jk}$ dans $\Sigma_k$ en \struck{relation de la forme} position relative correcte, telle que leur intersection soit la somme des $\Sigma_{ijk}$, on introduit comme plus haut une « fonction dimension » sur $I$ ; \struck{qui soit} \struck{$I_0=\ldots$}
\[
I\xrightarrow{\ \delta\ }\mathbb N
\]
qui soit strictement croissante, \struck{et on} \add{\ill{}} comme on l'a vu plus haut), \struck{et on construit}

\struck{(posant $X_{(d)}=\bigcup_{\delta(i)\le d}X_i$) $T_0$ voisinage tubulaire standard de $X_{(0)}$, $T_1$ voisinage tubulaire standard de $X_{(1)}\setminus\mathring T_0$ dans $X^{(1)}=X\setminus\dot X_{(0)}$}\note{Passage réuni par une accolade et traversé d'un long trait ondulé : lu comme biffé. Une formule en fin de ligne, sur $X_{(0)}$, est surchargée de hachures.}

On pose, pour $d\in\mathbb Z$
\[
(44)\qquad X_{(d)}=\bigcup_{\substack{i\in I\\ \delta(i)\le d}}X_i
\]
d'où
\[
(45)\qquad X_{-1}=\emptyset\subset X_{(0)}\subset X_{(1)}\subset\cdots\subset X_{(N)}=X\qquad\text{si } N=\sup_{i\in I}\delta(i)
\]
On construit par récurrence une filtration décroissante par des fermés
\[
(46)\qquad X^{(0)}=X\supset X^{(1)}\supset\cdots\supset X^{(N+1)}=\emptyset
\]
avec
\[
(47)\qquad X^{(i)}\subset X\setminus X_{i-1}
\]
i.e.\ $X^{(1)}\subset X\setminus X_{(0)}$, $X^{(2)}\subset X-X_{(1)}$, \ldots, $X^{(N+1)}\subset X\setminus X_{(N)}=X-X=\emptyset$.

\page{74}
\note{En haut à gauche, « 33 » de sa main.}

On choisit un voisinage tubulaire standard $T_{(0)}$ de $X_{(0)}$ dans $X$, et on pose
\[
(48)\quad\left\lbrace
\begin{array}{l}
X^{(1)}=X\setminus\mathring T_{(0)}\\[4pt]
X^{(1)}_{(d)}=X^{(1)}\cap X_{(d)}\qquad(=\emptyset\ \text{si}\ d\le0)\\[4pt]
X^{(1)}_i=X^{(1)}\cap X_i\qquad\text{de sorte que}\quad X^{(1)}_{(d)}=\bigcup_{\substack{i\in I\\ \delta(i)\le d}}X^{(1)}\cap X_i
\end{array}\right.
\]
On \struck{pose} \add{prend} alors un voisinage tubulaire \struck{standard de} standard $T_1$ de $X_{(1)}$ dans $X$ (petit par rapport au choix de $T_0$), soit $T^{(1)}_1$ sa trace dans $X^{(1)}$, qui est un vois.\ tub.\ standard de $X^{(1)}_{(1)}$ dans $X^{(1)}$. On pose
\[
(49)\qquad X^{(2)}=X^{(1)}\setminus\mathring T^{(1)}_{(1)}=X^{(1)}\cap(X\setminus\mathring T_1)=X\setminus(\mathring T_0\cup\mathring T_1)
\]
\note{Les deux numéros sont biffés en marge ; ils sont repris plus bas pour d'autres formules. Dans la dernière égalité, $X$ est écrit par-dessus une autre lettre.}

On prend un voisinage tubulaire $T_d$ de $X_{(d)}$ dans $X$, petit p.r.\ aux vois.\ tubulaires $T_{d'}$ ($d'<d$) déjà construits, et on pose
\[
(48)\qquad X^{(d)}=X-\bigcup_{\alpha<d}\mathring T_\alpha .
\]

\marginal{Ops $T_d=\bigcup_{\substack{i\in I\\ \delta(i)\le d}}T_i$, où $T_i$ vois.\ tub.\ de $X_i^{*}$ dans $X$}\note{en travers dans la marge gauche, en face de (48) ; « Ops » tel quel ; lecture incertaine.}

Ceci posé, on pose, pour $i\in I$, $d=\delta(i)$
\[
(49)\qquad
\begin{aligned}
\Sigma_i=X_i\cap X_{(d)}&=X_i\setminus\bigcup_{\alpha<d}X_i\cap\mathring T_\alpha\\
&=X_i\setminus\bigcup_{j<i}X_i\cap\mathring T_j\ \subset X_i^{*}
\end{aligned}
\]
\note{« $X_{(d)}$ » tel quel (on attend $X^{(d)}$). En fin de première ligne, un « $\cap X_i$ » est biffé ; sous « $j<i$ », un trait ondulé.}

\page{75}
\note{Sans numéro de sa main : suite du feuillet 33.}

Si $i<j$, \struck{$d=\delta(i)$, $d'=\delta(j)$}
\[
(49)\qquad \Sigma_{ij}=\dot T_i\cap\Sigma_j\subset\Sigma_j
\]
\note{Numéro (49) tel quel : il vient de servir, page 74. Après le signe $=$, un premier membre illisible est biffé.}

et pour un drapeau $i_0<i_1<\cdots<i_n$
\[
(50)\qquad
\begin{aligned}
\Sigma_{i_0\ldots i_n}&=\dot T_{i_0}\cap\dot T_{i_1}\cap\cdots\cap\dot T_{i_{n-1}}\cap\Sigma_{i_n}\\
&=\Sigma_{i_0i_n}\cap\Sigma_{i_1i_n}\cap\cdots\cap\Sigma_{i_{n-1}i_n}
\end{aligned}
\]
On a \add{de façon évidente} les morphismes de transition \add{entre les $\Sigma_{i_0\ldots i_n}$} pour les \struck{morphismes} inclusions \emph{cofinales} de drapeaux, mais non pour des inclusions quelconques, dont la définition dépend des choix de rétractions des
\[
T_i^{(d)}=T_i\cap X^{(d)}\qquad(d=\delta(i))
\]
sur $X_i^{(d)}=X_i\cap X^{(d)}$, ce qui donne alors
\[
\Sigma_{ij}\subset\dot T_i^{(d)}\subset T_i^{(d)}\longrightarrow X_i^{(d)}=\Sigma_i
\]
la structure complète $\bigl(\Sigma_i,\ (\Sigma_{ij}),\ \Sigma_{ij}\rightrightarrows\Sigma_j,\Sigma_i\bigr)$.

(Mais il faut des hyp.\ d'équisingularité sur la stratification des $X_i$, et une condition \add{de locale trivialité} \add{sur les} rétractions $T_i^{(d)}\to X_i$ i.e.\ $T^{(d)}_{(d)}\to X^{(d)}_{(d)}$, pour que les $\Sigma_{ij}\to\Sigma_i$ soient des fibrations)

\ldots

\section*{Introduction des provoisinages}

\page{76}
\note{En haut à gauche, « 34 » de sa main.}

15. \emph{Introduction des provoisinages} : $\mathcal V_{i_0\ldots i_n}$, $\mathcal V_J$,\note{Sur $\mathcal V_J$, et sur $\mathcal V_i$ dans les deux lignes qui suivent, un petit exposant biffé (une étoile ?).}

On désigne par $\mathcal V_i$ \struck{\ill{}} \add{pro}voisinage tubulaire (au sens des provoisinages) de $X_i^{*}=X_i\setminus\bigcup_{j<i}X_j$ dans $X$. Ainsi, \struck{\ill{}}
\[
(51)\qquad X_i^{*}\longrightarrow \mathcal V_i\qquad\text{équivalence d'homotopie des provoisinages.}
\]
\struck{Pour un drapeau} Si $i<j$, on pose
\[
(52)\qquad \mathcal V_{ij}=\mathcal V_i\cap \mathcal V_j
\]
de sorte qu'on a

\begin{tikzcd}
\mathcal V_{ij} \arrow[r, "b_{ij}"] \arrow[d, hook, "s_{ij}"'] & \mathcal V_j \\
\mathcal V_i &
\end{tikzcd}

$\bigl(X_i\xrightarrow[\approx]{(h)}\mathcal V_i\bigr)$

\marginal{\emph{NB} $\mathcal V_{ij}$ a un sens si $i\neq j$ mais si $i$ et $j$ ne sont pas comparables, on a $\mathcal V_i\cap \mathcal V_j=\emptyset$, si $i=j$ on a $\mathcal V_{ii}=\mathcal V_i$ (\ill{} si les ouverts de $X$ sont normaux, p.ex.\ $X$ métrisable)}\note{à droite, en face de (52) et du diagramme ; lecture incertaine dans le détail.}

\marginal{Pour un \ill{} $\mathcal V_{ij}^{\ill{}}=\mathcal V_i\cap$ \ldots\ équivalence d'homotopie ?}\note{en travers dans la marge gauche, en face de (52) ; lecture très incertaine.}

On a :
\[
(53)\qquad X=\varinjlim\bigl(\mathcal V_i,\ \mathcal V_{ij},\ \rightrightarrows b_{ij},s_{ij}\bigr)
\]
$X$ est obtenu par \emph{recollement} des $\mathcal V_i$ suivant les $\mathcal V_{ij}$.

On définit aussi, plus généralement, les
\[
(54)\qquad \mathcal V_{i_0\ldots i_n}=\mathcal V_{i_0}\cap\cdots\cap \mathcal V_{i_n}
\]
pour un drapeau $i_0<i_1<\cdots<i_n$

\marginal{\emph{NB} L'expression a un sens pour $(i_0,\ldots,i_n)$ quelconque, mais \ill{} vide si ce n'est pas un drapeau}\note{en travers dans la marge droite, en face de (54) ; lecture incertaine.}

\page{77}
\note{Sans numéro de sa main : suite du feuillet 34.}

c'est un système inductif \ill{} provoisinages sur l'ens.\ ordonné opposé de celui des drapeaux de $I$ (avec l'inclusion), dont la \emph{limite} est égale $\approx X$.

Soit $Z$ \struck{\ill{}} une partie (\add{loc\textsuperscript{t}} fermée) de $X$, et posons, pour \struck{\ill{}} un $i$ tel que \struck{$X_i\subset X$}
\[
(55)\qquad \mathcal V_i^{Z}=\mathcal V_i\cap Z\qquad\text{provoisinage tubulaire de } X_i^{*} \text{ dans } Z
\]
\struck{et} on a
\[
(56)\qquad X_i^{*}\hookrightarrow \mathcal V_i^{Z}\hookrightarrow \mathcal V_i\qquad\text{équivalences d'homotopie}
\]
Plus généralement, on pose (si $X_{i_n}\subset Z$)
\[
\mathcal V^{Z}_{i_0\ldots i_n}=\mathcal V_{i_0\ldots i_n}\cap Z=\mathcal V^{Z}_{i_0}\cap\cdots\cap \mathcal V^{Z}_{i_n}
\]
C'est aussi le $\mathcal V_{i_0\ldots i_n}$ relatif : la \struck{\ill{}} stratification de $Z$ induite par les $Z_i\overset{\text{déf}}{=}X_i\cap Z$, pour le système d'indices $i_0\ldots i_n$. On a
\[
(57)\qquad \mathcal V^{Z}_{i_0\ldots i_n}\hookrightarrow \mathcal V_{i_0\ldots i_n}\qquad\text{équiv.\ d'homotopie}
\]
Soit $J$ une partie de $I$, posons
\[
(58)\qquad X_J^{*}=\bigcup_{i\in J}X_i^{*},\qquad X_J=\bigcup_{i\in J}X_i
\]

\page{78}
\note{En haut à gauche, « 35 » de sa main.}

Si $J$ est un « idéal » i.e.\ $i\in J$, $j<i\Rightarrow j\in J$, alors on a $X_J^{*}=X_J$, donc $X_J^{*}$ est fermé.

Si $J$ est une « bande » i.e.\ si $i,j\in J$, $i<j$, alors $\forall k\in I$, $(i<k<j\Rightarrow k\in J)$, ou ce qui revient au même, si $J'=\widetilde J\setminus J$ est un « idéal », où $\widetilde J$ est l'idéal engendré par $J$, alors on a
\[
(59)\qquad X_J^{*}=X_{\widetilde J}\setminus X_{J'}
\]
différence de deux parties fermées, c'est une partie loc\textsuperscript{t} fermée de $X$, dont les
\[
X_i^{J}\overset{\text{déf}}{=}X_i\cap X_J^{*}\qquad(i\in J)
\]
forment une stratification.\note{Sous (59), une formule biffée.} Posons
\[
(60)\qquad \mathcal V^{J}_{j_0\ldots j_n}=\mathcal V_{j_0\ldots j_n}\cap X_J^{*}=\mathcal V^{Z}_{j_0\ldots j_n}\qquad\text{pour } Z=X_J^{*}
\]
pour $j_0,\ldots,j_n\in J$, $j_0<j_1<\cdots<j_n$, \struck{\ill{}} ce sont alors les $\mathcal V_{i_0\ldots i_n}$ associés à la stratification de type $J$ de $X_J^{*}=Z$. On a donc
\[
(61)\qquad
\begin{aligned}
X_J^{*}&\simeq\varinjlim_{i,j\in J}\bigl(\mathcal V_i^{J},\mathcal V_{ij}^{J}\bigr)\\
&\simeq\varinjlim_{(i_0<\cdots<i_n)\in\mathrm{Drap}(J)}\bigl(\mathcal V^{J}_{i_0\ldots i_n}\bigr)
\end{aligned}
\]
\struck{On s'intéresse : la situation des}

\page{79}
\note{Sans numéro de sa main : suite du feuillet 35. Toute la page est barrée de deux longues diagonales ; elle se lit néanmoins, et elle est donnée entière. La page 80 reprend ses numéros (62) et (63).}

deux bandes $J\subset J'$, d'où
\[
X_J^{*}\subset X_{J'}^{*}
\]
et on prendra le \add{pro}voisinage tubulaire de $X_J^{*}$ dans $X_{J'}^{*}$, on trouve
\[
\mathcal V_{X_J^{*},X_{J'}^{*}}=\varinjlim_{i,j\in J}\bigl(\mathcal V_i^{J'},\mathcal V_{ij}^{J'}\bigr)\qquad\Bigl(X_J^{*}=\varinjlim_{i,j\in J}\bigl(\mathcal V_i^{J},\mathcal V_{ij}^{J}\bigr)\Bigr)
\]
\note{La formule entre parenthèses est entourée ; une flèche la relie, par « (62) », au premier membre. Un mot biffé sous « $i,j\in J$ ».}

\add{de plus} \uncertain{Ops} $J'=I$, \add{$J$ un} idéal dans $J'=I$, (\emph{NB} quitte à remplacer $I$ par un intervalle, $J_0=\widetilde J\setminus J$ ; \struck{\ill{}} \add{on est ramené au cas où} $I\supset J_0\subset J_1\subset J_2$ trois idéaux\ldots\ $J=J_1\setminus J_0$, $J'=J_2\setminus J_0$ \ldots).

Soit maintenant, de plus, $K\subset I$ une partie de $I$, et considérons \struck{\ill{}} $\mathcal V_J^{J'}\setminus X_K\cap \mathcal V_J^{J'}=$
\[
\mathcal V^{J'}_{J;K}=\mathcal V(X_J^{*},X_{J'}^{*})\setminus X_K\cap \mathcal V(X_J^{*},X_{J'}^{*})
\]
Notons que pour $k\in I$
\[
(63)\qquad \mathcal V^{J'}_{J;k}\overset{\text{déf}}{=}\mathcal V^{J'}_{J}\cap X_k=\emptyset\ \Longleftrightarrow\ X_J^{*}\cap X_k=\emptyset
\]
\struck{$\Longleftrightarrow\ \forall i\in I,\ i\le k\Rightarrow i\notin J$}

et comme $X_k$ est la réunion (disjointe) des $X_i^{*}$ pour $i\le k$, la relation (63) est conséquence de celle que $i\le k\Rightarrow i\notin J$ i.e.\ \struck{\ill{}}

\marginal{On pourrait \ill{} (62) le $\mathcal V_J^{J'}$. On $\mathcal V_J^{J'}=\mathcal V_J^{J'_0}$ \ldots\ $J_0\subset J'$ \ldots\ idéal \ldots\ la condition $X_J^{*}$ \ldots\ fermé dans $X_{J'}^{*}$ \ldots}\note{notes en travers dans la marge gauche, sur toute la hauteur ; lecture très incertaine.}

\page{80}
\note{En haut à gauche, « 36 » de sa main.}

\struck{Le cas des $\mathcal V$}

Posons, pour une bande $J$
\[
(62)\qquad
\begin{aligned}
\mathcal V_J&=\text{provoisinage tubulaire de } X_J^{*} \text{ dans } X\\
&=\bigcup_{i\in J}\mathcal V_i=\varinjlim_{i,j\in J}\bigl(\mathcal V_i,\mathcal V_{ij}\bigr)
\end{aligned}
\]
\note{Le numéro (62), surchargé, reprend celui de la page 79, barrée.}

On s'intéresse aux provoisinages de la forme
\[
(63)\qquad \mathcal V_J^{L}\overset{\text{déf}}{=}\mathcal V_J\cap X_L^{*},
\]
où $J$, $L$ sont deux bandes. Par exemple, on retrouve les $X_J^{*}$ comme
\[
(64)\qquad \mathcal V_J^{J}=X_J^{*}
\]

\marginal{\emph{Ex} $J\subset\hat J'$ inclusion de bandes, d'où $X_J^{*}\subset X_{J'}^{*}$, on prend $\mathcal V(X_J^{*},X_{J'}^{*})\setminus \mathcal V(X_J^{*},X_{J'}^{*})\cap X_K^{*}$ [$K$ (idéal de $J'$)]}\note{à droite de (63), relié par un trait à « la forme ».}

\marginal{provoisinage de \ill{} second ordre entre $X_J^{*}$ et $X_L^{*}$ : $\mathcal V_J^{L}\hookrightarrow X_L^{*}\hookrightarrow \mathcal V_L$ \ldots}\note{en travers dans la marge gauche, en face de (63)-(64), suivie de griffonnages biffés ; lecture incertaine.}

Soit $L_0$ l'ens.\ des $l\in L$ qui majorent un élément de $J$, et $J_0$ l'ens.\ des $j\in J$ qui minorent un élément de $L$, on a alors
\[
(65)\qquad \mathcal V_J^{L}=\mathcal V_J^{L_0}=\mathcal V_{J_0}^{L}=\mathcal V_{J_0}^{L_0}
\]
i.e.\ $J_0$, $L_0$ sont des sous-bandes de $J$, $L$ respectivement ($J_0$ est un idéal de $J$, $L_0$ un coidéal de $L$) telles que tout élément de $J_0$ soit majoré par un élément de $L_0$, tout élément de $L_0$ majore un élément de $J_0$. On est ainsi ramené au cas de bandes dans la relation de domination
\[
(66)\qquad J\prec L
\]
qui signifie que tout élément de $J$ est majoré par un élément de $L$, tout élément de $L$ majore un élément de $J$.

donc

\end{document}
